\documentclass[12pt,a4paper]{article}
\usepackage{ugmath}
\usepackage{float}
\usepackage{placeins}
\usepackage[labelfont=bf,labelsep=space]{caption}
\newcommand{\paren}[1]{\left( #1 \right)}
\newcommand{\alumno}{Ricardo León Martínez}
\newcommand{\materia}{Fisica I}
\newcommand{\profesor}{Lucero Uscanga Aguilera}
\newcommand{\tarea}{Tarea 4}
\newcommand{\fecha}{13/3/2026}

\begin{document}
    \begin{center}
    {\large\textbf{UNIVERSIDAD DE GUANAJUATO}}\\[0.3cm]
    {\normalsize\textbf{DIVISIÓN DE CIENCIAS NATURALES Y EXACTAS}}\\
    {\normalsize\textbf{CAMPUS GUANAJUATO}}\\[1cm]

    {\Large\textbf{\tarea\ (\materia)}}\\[1cm]
    \end{center}

    \textbf{Nombre:} \alumno \hfill 
    \textbf{Fecha:} \fecha \hfill 
    \textbf{Calificación:} \rule{3cm}{0.4pt} \\[0.3cm]

    \begin{mdframed}[style=mdbluebox,frametitle={Ejercicio 1}]
        Un cuerpo, inicialmente en reposo $(\theta=0$ y $\omega=0$ cuando $t=0)$ es
        acelerado en una trayectoria circular de $1.3$m de radio, de acuerdo a la ecuacion
        $\alpha=120t^{2}-48t+16$. Encontrar la posición angular y la velocidad angular del
        cuerpo en función del tiempo, y las componentes tangencial y centrípeta de la
        aceleración.
    \end{mdframed}

    \begin{mdframed}[style=mdbluebox,frametitle={Ejercicio 1}]
        Una pequeña masa $m$ se mueve sobre la mesa horizontal sin fricción. Está unida a
        una cuerda ligera que pasa por un pequeño agujero en la mesa. Del otro lado de la cuerda,
        inicialmente la partícula se mueve en una circunferencia de radio $r_{0}$ con velocidad
        tangencial $v_{0}$. Luego la cuerda se acorta y la partícula queda moviéndose en una
        circunferencia de radio $r$.
        \begin{enumerate}
            \item[(a)] Calcular el momento angular inicial de la partícula respecto al agujero.
            \item[(b)] Usar conservación del momento angular para encontrar la nueva velocidad tangencial $v(r)$
            \item[(c)] Determinar cómo cambia la energía cinética de la partícula.
            \item[(d)] Explicar por qué la energía cinética cambia aunque el momento angular se conserve.   
        \end{enumerate}
    \end{mdframed}
\end{document}